%% Generated by scripts/export_prompts_appendix.py. Do not edit by hand.

\onecolumn
\sloppy

\lstset{basicstyle=\ttfamily\scriptsize, breaklines=true, breakatwhitespace=false, columns=fullflexible, keepspaces=true, showstringspaces=false, frame=single, framerule=0.3pt, aboveskip=4pt, belowskip=4pt, xleftmargin=3pt, xrightmargin=3pt, upquote=true}

\section{Prompts}\label{app:prompts}
Every prompt is printed as the experiment code builds it. Fields filled at run time appear as \texttt{<{}<NAME>{}>}, and non-ASCII arrows and dashes are printed in ASCII. Gemini 2.5 Flash answers every call unless stated otherwise, and the Gemini 2.5 Pro runs of Appendix~\ref{app:pro} use the same prompts.

\subsection{VerifiQuant at run time}
The model is called in exactly three places. Each call requests JSON through a response schema, whose fields are listed after the prompt.

\paragraph{Contract selection (N and M)} One block per retrieved candidate, top 4, in rank order.
\begin{lstlisting}
You are selecting a financial formula card from candidates.

Decision policy -- judge ONLY the user's INTENT (which metric they want), NOT the quality of the
input values. Apply in order:
1) decision="select_card" -- if the QUESTION or the CONTEXT explicitly names a specific
   metric/formula (e.g. "NPV", "Fixed Asset Turnover ratio", "monthly loan payment", "CAGR",
   "working ratio") and a candidate matches that named metric. This INCLUDES a clarification the
   user added to the context (e.g. a note like "use NPV" or "the answer is: Net Present Value") --
   once the user has named the metric, the ambiguity is resolved, so SELECT that card; do NOT
   abstain_m again. SELECT IT even if the provided numbers look wrong, negative, out of range,
   missing, or otherwise anomalous -- validating input values is handled by later deterministic F/E
   checks, NOT by you. Do not abstain because the data looks problematic.
2) decision="abstain_m" -- ONLY when the user states a GENERAL GOAL without naming a metric
   (e.g. "is this project worth it", "will it earn money", "how did this investment perform") AND
   two or more candidates could each plausibly serve that goal. You MUST abstain_m here even if the
   variables perfectly fit one candidate; DO NOT guess the "most appropriate"/"best"/"industry-
   standard" metric or argue others are "less suitable" -- that choice belongs to the user. Provide
   1-3 clarification_questions naming the competing metrics (e.g. "NPV, discounted payback, or ROI?").
3) decision="abstain_n" -- the named metric/logic is genuinely NOT supported by ANY candidate (the
   library has no card for it), OR the user asks you to invent/derive a new formula. Provide
   support_gap_reason. A supported metric with bad input values is NOT abstain_n -- select it.

Reminder: a named metric (even with bad data) -> select_card. A vague goal with several fitting
metrics -> abstain_m. Truly unsupported/invented metric -> abstain_n.

Question:
<<QUESTION>>

Context:
<<CONTEXT>>

Candidates:
Candidate <<RANK>> | fic_id=<<FIC_ID>> | score=<<RETRIEVAL SCORE>>
title=<<TITLE>>
topic=<<TOPIC>>
summary=<<SUMMARY>>
selection_hints=<<SELECTION_HINTS>>
applicable_when=<<APPLICABLE_WHEN>>
not_applicable_when=<<NOT_APPLICABLE_WHEN>>
scope_boundaries=<<SCOPE_BOUNDARIES>>

Return JSON only.
\end{lstlisting}
Output fields: \texttt{decision} (select\_card, abstain\_m, abstain\_n), \texttt{chosen\_fic\_id}, \texttt{reason}, \texttt{support\_gap\_reason}, \texttt{ambiguity\_tags}, \texttt{clarification\_questions}.

\paragraph{Input binding (F and E read the bound values)}\leavevmode
\begin{lstlisting}
You are extracting structured input values for a financial formula.

For each required input, locate its value in the Question / Context and:
1. Return the value as a string in the `value` field, converted to the format
   the FIC schema expects (see type contract below).
2. Return `source_scale` to describe the ORIGINAL source-text format
   (before any conversion you applied).

Type contract for `value`:
- type "number" / "decimal" / "float": numeric string, e.g. "0.08".
- type "integer": whole-number string, e.g. "5".
- type "boolean": "true" or "false".
- type "array[number]": flat JSON array string, e.g. "[100, 200, 300]".
- type "array[array[number]]" or nested array types:
    JSON array-of-arrays, e.g. "[[50, 240], [80, 220]]".
    Do NOT flatten nested arrays into a single list.
- type "dict[string,number]": JSON object string with string keys and numeric
  values, e.g. "{"Sales": 3000, "R&D": 5000}".
- type "array[object]": JSON array of objects, e.g.
  "[{"quantity": 100, "price": 50}, {"quantity": 50, "price": 48}]".
- type "object": JSON object or array string matching the input description.
- unit "decimal" or "decimal_rate": convert % to decimal (15% -> 0.15).
- unit "%" or "percent": keep as percentage points (15% -> 15).

Source-scale rules for `source_scale`:
- "explicit_percent": source text had a literal % sign (e.g. "15%", "8 percent").
- "explicit_decimal": source text was already in decimal form (e.g. "0.08", "0.15").
- "explicit_integer": source text was an unambiguous integer with no % sign
  (e.g. "5 years", "120 units").
- "inferred": you had to infer or assume the scale from context -- be honest.

If a value is absent, set status="missing", value="", source_scale="inferred".
Document every normalisation decision in normalization_note.

Required inputs (full schema):
<<CONTRACT INPUT SCHEMA (JSON: name, type, unit, description, aliases)>>

Question:
<<QUESTION>>

Context:
<<CONTEXT>>

Return JSON only.
\end{lstlisting}
Output fields: \texttt{inputs}, \texttt{normalization\_note}, where each input carries \texttt{name}, \texttt{status}, \texttt{value}, and \texttt{source\_scale}.

\paragraph{Ambiguity critic (I)} When the contract declares no semantic hints, the two mode rules are replaced by the open-mode rules shown second.
\begin{lstlisting}
You are the VerifiQuant I-gate Critic Agent.
Goal: detect hidden semantic ambiguity before deterministic execution.

Ambiguity examples:
- FX direction (USD/TWD vs TWD/USD)
- time basis (begin-year vs end-year cash flow)
- compounding convention (monthly vs annual)

Rules:
- If ambiguity materially changes calculation interpretation, set needs_clarification=true.
- Provide concise clarification questions.
- Mode: hint_oriented
- semantic_hints are available: use them as primary reference and fill triggered_hint_ids when relevant.
- If a hint is not relevant, do not trigger it.


Question:
<<QUESTION>>

Context:
<<CONTEXT>>

Semantic hints:
<<DECLARED SEMANTIC HINTS (JSON)>>

Provided inputs:
<<BOUND INPUTS (JSON)>>

Return JSON only.
\end{lstlisting}
\begin{lstlisting}
- semantic_hints are empty: perform OPEN semantic ambiguity scan from question/context.
- In open mode, triggered_hint_ids should be empty.
- If clarification is needed, propose concise clarification_options when possible.
\end{lstlisting}
Output fields: \texttt{needs\_clarification}, \texttt{triggered\_hint\_ids}, \texttt{ambiguity\_tags}, \texttt{clarification\_questions}, \texttt{clarification\_options}, \texttt{reason}.

\subsection{Clarification oracle ($K{>}1$)}
The oracle sees the reference solution code of the question, never the gold number, and may only rewrite the question and context. VerifiQuant's oracle reads the DiagnosticReport of a structured exit. The CoT oracle of Table~\ref{tab:v2kabl} reads the CoT output of every turn.

\paragraph{VerifiQuant oracle.}\leavevmode
\begin{lstlisting}
You are an Oracle-in-the-loop rewriter for VerifiQuant framework-guided iteration.
You can ONLY use the logic within the ground-truth code to clarify assumptions and revise user question/context. You must NOT rely on the final numeric ground truth result.

Goal:
- Given framework diagnostic output, rewrite question/context so the next run can pass gates.
- Keep semantics faithful to original task and the logic implied by the provided code.
- Do NOT output the final numeric answer directly in question/context.

Current question:
<<QUESTION>>

Current context:
<<CONTEXT>>

Framework diagnostic:
<<DIAGNOSTICREPORT OF THE PREVIOUS TURN (JSON)>>

Ground-truth code:
<<REFERENCE SOLUTION CODE>>

Return JSON only.
\end{lstlisting}
Output fields: \texttt{updated\_question}, \texttt{updated\_context}, \texttt{notes}.

\paragraph{CoT oracle.}\leavevmode
\begin{lstlisting}
You are Oracle-in-the-loop support for a pure CoT baseline.
There is no VerifiQuant framework in this loop.
You can ONLY use the logic within the ground-truth code to clarify assumptions and revise user question/context. You must NOT rely on the final numeric ground truth result.

Current question:
<<QUESTION>>

Current context:
<<CONTEXT>>

Current CoT step output:
<<COT OUTPUT OF THE PREVIOUS TURN (JSON)>>

Ground-truth code:
<<REFERENCE SOLUTION CODE>>

Return JSON only.
\end{lstlisting}
Output fields: \texttt{updated\_question}, \texttt{updated\_context}, \texttt{notes}.

\subsection{Baselines}
R-PAL and R-PAL$^{+}$ receive one reference block per retrieved candidate, top 4, in rank order. The same R-PAL prompt serves the clean questions, the traps, the out-of-library split, and the human-labeled set, and the same holds for R-PAL$^{+}$.

\paragraph{CoT (single shot)}\leavevmode
\begin{lstlisting}
You are a financial CoT solver with self-improvement hints.
Solve using the current question/context and return a tentative answer.
If information is missing/ambiguous, set needs_more_info=true and list missing items.
IMPORTANT: `answer` must be a numeric string only (for example: "10185.19"), or empty string if you cannot provide one.

Question:
<<QUESTION>>

Context:
<<CONTEXT>>

Return JSON only.
\end{lstlisting}
Output fields: \texttt{answer}, \texttt{needs\_more\_info}, \texttt{missing\_information}, \texttt{confidence}, \texttt{revised\_question}, \texttt{revised\_context}, \texttt{reasoning\_note}.

\paragraph{PAL (no retrieval)}\leavevmode
\begin{lstlisting}
You are a financial reasoning assistant that solves problems by writing Python.
Write a single self-contained Python function solve() that computes the numeric answer.
Rules:
- Only the Python standard library `math` may be imported; no other imports.
- No input(), no file or network access.
- solve() must return a single number (int or float).
- If the problem cannot be solved from the given information, return the string "CANNOT SOLVE".
Return ONLY the code, in one ```python block.

Question:
<<QUESTION>>

Context:
<<CONTEXT>>
\end{lstlisting}

\paragraph{R-PAL.}\leavevmode
\begin{lstlisting}
You are a financial reasoning assistant that solves problems by writing Python.

You are given top-ranked reference functions retrieved from a financial function library.
The candidates are retrieval results, not guaranteed to be relevant or correct for the question.
Select the useful candidate, if any, and adapt its logic to the exact question and context.

Write a single self-contained Python function solve() that computes the numeric answer.
Rules:
- Only the Python standard library `math` may be imported; no other imports.
- No input(), no file or network access.
- solve() must return a single number (int or float).
- Match the output unit and scale requested by the question.
- If no candidate applies or the problem cannot be solved from the given information,
  return the string "CANNOT SOLVE".
- Do not merely call a reference function: solve() must be self-contained.
- Return ONLY the code, in one ```python block.

Question:
<<QUESTION>>

Context:
<<CONTEXT>>

Retrieved reference functions (ranked):
### Candidate <<RANK>>: <<TITLE>>
Summary: <<SUMMARY>>
```python
<<REFERENCE FUNCTION CODE>>
```
\end{lstlisting}

\paragraph{R-PAL$^{+}$ (taxonomy prompt)}\leavevmode
\begin{lstlisting}
You are a financial reasoning assistant using retrieved reference functions.

Before calculating, decide whether the request is safe and sufficiently specified.
Use this taxonomy and priority order:
- M: the user's intent or requested metric is unclear; ask what they want to evaluate.
- N: the requested calculation is unsupported or the references are only superficially related; refuse.
- F: a required numeric/input field is absent or redacted; ask for the exact missing field.
- E: a provided value violates a financial, mathematical, or domain boundary; identify it and do not calculate.
- I_HARD: an unresolved unit, direction, timing, basis, or convention can change the answer; clarify before calculating.
- I_SOFT: the calculation can proceed, but an interpretation assumption must be disclosed as a warning.
- READY: the request has one clear intent, supported logic, complete inputs, admissible values, and no material ambiguity.

The retrieved candidates are not guaranteed to be applicable or safe. Do not force-fit one.
For READY or I_SOFT, write a self-contained Python solve() function. Only `math` may be imported.
For M/N/F/E/I_HARD, do not provide executable code or a numeric answer.

Return ONLY one valid JSON object with exactly these fields:
{
  "decision": "execute|warn|clarify|refuse",
  "diagnostic": "READY|M|N|F|E|I_HARD|I_SOFT",
  "reason": "concise evidence-grounded explanation",
  "action_targets": ["exact missing field, invalid field/rule, or ambiguity target"],
  "clarification_question": "one actionable question or empty string",
  "warning": "explicit assumption/warning for I_SOFT or empty string",
  "code": "self-contained Python code as a JSON string, or empty string"
}

Question:
<<QUESTION>>

Context:
<<CONTEXT>>

Retrieved reference functions (ranked):
### Candidate <<RANK>>: <<TITLE>>
Summary: <<SUMMARY>>
```python
<<REFERENCE FUNCTION CODE>>
```
\end{lstlisting}

\paragraph{Self-consistency.} Five samples at temperature 1.0, majority vote over answers within 1\%.
\begin{lstlisting}
You are a financial reasoning assistant. Think step by step, then give the final numeric answer.
End with exactly one line: 'Final answer: <number>'.

Question:
<<QUESTION>>

Context:
<<CONTEXT>>
\end{lstlisting}

\subsection{Staged controls}
Every stage prompt starts with the taxonomy block of the R-PAL$^{+}$ prompt above, verbatim from its first line through ``Do not force-fit one.'', shown here as \texttt{<{}<TAXONOMY>{}>}. Lines marked \texttt{[Staged+Spec only:]} are absent from Staged.

\paragraph{Stage 1: scope.}\leavevmode
\begin{lstlisting}
<<TAXONOMY>>

The decision is made in five separate stages, in this order:
(1) scope: M and N; (2) inputs: F; (3) boundaries: E; (4) ambiguity: I_HARD and I_SOFT; (5) program.
Any earlier stages have passed. This is stage 1, scope. Decide only whether the retrieved candidates support the requested calculation. If at least one does, choose the most applicable candidate. Return M if the user's intent or requested metric is unclear, or N if the requested calculation is unsupported or the references are only superficially related.
Later stages handle the other categories, so do not judge them here.

Return ONLY one valid JSON object with exactly these fields:
{
  "diagnostic": "PASS|M|N",
  "chosen_candidate": <candidate number, or 0 for M/N>,
  "reason": "concise evidence-grounded explanation",
  "clarification_question": "one actionable question for M, or empty string"
}

Question:
<<QUESTION>>

Context:
<<CONTEXT>>

Retrieved reference functions (ranked):
### Candidate <<RANK>>: <<TITLE>>
Summary: <<SUMMARY>>
```python
<<REFERENCE FUNCTION CODE>>
```  [repeated for each candidate]
[Staged+Spec only:]
Contract metadata:
selection_hints: ["<<SELECTION_HINTS>>"]
applicable_when: ["<<APPLICABLE_WHEN>>"]
not_applicable_when: ["<<NOT_APPLICABLE_WHEN>>"]
scope_boundaries: ["<<SCOPE_BOUNDARIES>>"]
\end{lstlisting}

\paragraph{Stage 2: inputs.}\leavevmode
\begin{lstlisting}
<<TAXONOMY>>

The decision is made in five separate stages, in this order:
(1) scope: M and N; (2) inputs: F; (3) boundaries: E; (4) ambiguity: I_HARD and I_SOFT; (5) program.
Any earlier stages have passed. This is stage 2, inputs. Decide only whether every input the calculation needs is present. Use the chosen reference function's parameters as the list of inputs; a parameter with a default value is not missing. Bind each input to its value in the question or context. Return F if a required numeric/input field is absent or redacted.
Later stages handle the other categories, so do not judge them here.

Return ONLY one valid JSON object with exactly these fields:
{
  "diagnostic": "PASS|F",
  "inputs": [{"name": "parameter name", "status": "provided|missing", "value": "value as stated, or empty string"}],
  "missing_fields": ["exact missing field"],
  "reason": "concise evidence-grounded explanation",
  "clarification_question": "one actionable question for F, or empty string"
}

Question:
<<QUESTION>>

Context:
<<CONTEXT>>

<<CHOSEN REFERENCE BLOCK>>
[Staged+Spec only:]

Contract input schema:
<<CONTRACT INPUT SCHEMA (JSON: name, type, required, unit, description)>>
\end{lstlisting}

\paragraph{Stage 3: boundaries.}\leavevmode
\begin{lstlisting}
<<TAXONOMY>>

The decision is made in five separate stages, in this order:
(1) scope: M and N; (2) inputs: F; (3) boundaries: E; (4) ambiguity: I_HARD and I_SOFT; (5) program.
Any earlier stages have passed. This is stage 3, boundaries. Decide only whether a provided value violates a financial, mathematical, or domain boundary. Return E if one does.
Later stages handle the other categories, so do not judge them here.

Return ONLY one valid JSON object with exactly these fields:
{
  "diagnostic": "PASS|E",
  "invalid_fields": ["invalid field or violated rule"],
  "reason": "concise evidence-grounded explanation",
  "clarification_question": "one actionable question for E, or empty string"
}

Question:
<<QUESTION>>

Context:
<<CONTEXT>>

<<CHOSEN REFERENCE BLOCK>>

Bound inputs (from stage 2):
<<BOUND INPUTS FROM STAGE 2 (JSON)>>
[Staged+Spec only:]

Contract boundary checks (declared rules; judge them, they are not run):
<<DECLARED E-CHECKS (JSON: rule_id, name, description, expression, applies_to)>>
\end{lstlisting}

\paragraph{Stage 4: ambiguity.}\leavevmode
\begin{lstlisting}
<<TAXONOMY>>

The decision is made in five separate stages, in this order:
(1) scope: M and N; (2) inputs: F; (3) boundaries: E; (4) ambiguity: I_HARD and I_SOFT; (5) program.
Any earlier stages have passed. This is stage 4, ambiguity. Decide only whether an interpretation question remains. Return I_HARD if an unresolved unit, direction, timing, basis, or convention can change the answer, I_SOFT if the calculation can proceed but an interpretation assumption must be disclosed as a warning, or PASS if neither.
Later stages handle the other categories, so do not judge them here.

Return ONLY one valid JSON object with exactly these fields:
{
  "diagnostic": "PASS|I_HARD|I_SOFT",
  "ambiguity_targets": ["ambiguity target"],
  "reason": "concise evidence-grounded explanation",
  "clarification_question": "one actionable question for I_HARD, or empty string",
  "warning": "explicit assumption/warning for I_SOFT, or empty string"
}

Question:
<<QUESTION>>

Context:
<<CONTEXT>>

<<CHOSEN REFERENCE BLOCK>>

Bound inputs (from stage 2):
<<BOUND INPUTS FROM STAGE 2 (JSON)>>
[Staged+Spec only:]

Contract semantic hints (declared ambiguities):
<<DECLARED SEMANTIC HINTS (JSON)>>
\end{lstlisting}

\paragraph{Stage 5: program.}\leavevmode
\begin{lstlisting}
<<TAXONOMY>>

The decision is made in five separate stages, in this order:
(1) scope: M and N; (2) inputs: F; (3) boundaries: E; (4) ambiguity: I_HARD and I_SOFT; (5) program.
Any earlier stages have passed. This is stage 5, program. Write a self-contained Python solve() function that computes the answer from the bound inputs. Only `math` may be imported. If an I_SOFT assumption is listed below, follow it.
Later stages handle the other categories, so do not judge them here.

Return ONLY one valid JSON object with exactly these fields:
{
  "code": "self-contained Python code as a JSON string"
}

Question:
<<QUESTION>>

Context:
<<CONTEXT>>

<<CHOSEN REFERENCE BLOCK>>

Bound inputs (from stage 2):
<<BOUND INPUTS FROM STAGE 2 (JSON)>>
[only if stage 4 returned I_SOFT:]

I_SOFT assumption (from stage 4):
<<WARNING FROM STAGE 4>>
\end{lstlisting}

\subsection{Library construction}
Each formula family is built by three model calls. The first writes the core card from the source record, which holds the family's reference function together with the question it was sampled from, that question's context, and its worked solution (\texttt{python\_solution}), but never its numeric gold answer. The second writes retrieval metadata and the third writes repair rules, both from the core card. Deterministic normalization, lint, and the type audit of \S\ref{sec:v2method} follow. 176 of the 250 cards were built on June 15, 2026 with the core prompt below, as committed on May 23. The other 74 were carried over from earlier builds made with earlier revisions of the same prompt, all kept in the repository history. A July revision added an input and output type policy, which the library meets through the build-time type audit rather than regeneration.

\paragraph{Question details in cards.} Because the first call sees the source question, a card can repeat details of it. We count a card as doing so when one of its fields holds a number that appears in its source question or worked solution but not in its reference function, ignoring common constants (0 to 5, 10, 12, 24, 30, 52, 60, 100, 252, 360, 365, 1000, and a few unit values). 86 of the 250 cards do: 68 in a declared semantic hint, 26 in a boundary check, 22 in an input description, and 1 in the compute code. Most are example values inside a declared ambiguity, such as a hint that asks whether a stated rate of 7\% is meant as 0.07 or as 7.

\paragraph{Core card (stage 1)}\leavevmode
\begin{lstlisting}
Role:
You are a financial software engineer generating a FIC core card.

Task:
Convert one dataset case into a reusable `fic_core` object with:
- fic_id, name, short_description
- domain, topic
- inputs, output
- execution (deterministic python)
- selection_hints
- semantic_hints
- diagnostic_checks

Critical rules:
1) `python_solution` is PRIMARY for executable logic.
2) `function` is SECONDARY for naming and financial semantics.
3) Do not hardcode case-specific constants into execution; use inputs.
4) Use `article_content_excerpt` as supporting context to generate higher-quality `semantic_hints`.
5) If a diagnostic rule expression needs to validate the formula's output, you MUST call `compute(inputs)` to obtain the output value. You CANNOT access output variables directly from the `inputs` dictionary.

Taxonomy policy:
- Domain MUST be one of taxonomy domains.
- Topic should prefer existing topics under that domain.
- If no existing topic is suitable, you MAY create a concise snake_case topic under that existing domain.
- Never invent a new domain.

Diagnostic rules policy:
- Provide 3 to 8 checks.
- check_type must be one of: deterministic, normalization.
- deterministic checks must include a valid python boolean `expression` using input names. If the expression needs to check the computed output, you MUST call `compute(inputs)` instead of reading the output name from the `inputs` dictionary.
- For deterministic/normalization checks, include `predicate_mode`:
  - `violation`: expression=True means the rule is triggered (preferred default).
  - `validity`: expression=True means the input is valid, and expression=False triggers the rule.
- In diagnostics checks, use ONLY F or E. Do NOT output M/N/I/C.
- Why: M is handled at retrieval/card-selection stage before card commitment.
- Mapping guidance:
  - F is STRICTLY for schema/spec incompleteness only:
    - required input missing
    - input cannot be parsed to expected type
    - malformed structure (for example expected list but got scalar and cannot normalize)
    - for array inputs whose elements are tuples/pairs (e.g. each element must be
      [quantity, price]), the STRUCTURE check (each element is a list/tuple of the
      expected length) is F -- because malformed structure means the input cannot be
      parsed to the expected schema, even if numbers were provided.
  - E is for value-level or interpretation-level issues AFTER inputs exist:
    - unit/scale mismatch (8 vs 0.08, 15 vs 0.15)
    - bound/plausibility violations (negative rate, impossible ratio)
    - binding/direction/frequency inconsistencies
  - Prefer E over F when a value is present but suspicious or needs confirmation.
  - Do NOT classify percentage-range checks as F when the value exists; classify them as E with deterministic alert.
  - "Hard formula preconditions" should generally be E unless the variable is truly missing/unparseable.
  - If an array input is described as a list of (quantity, price) pairs or similar
    multi-element tuples, you MUST include an F-class check that verifies the structure
    of each element (e.g. `all(isinstance(item, (list, tuple)) and len(item) == 2 for item in inputs['field'])`)
    with predicate_mode="validity".
- severity must be error or alert.
- semantic_hints are for I-gate only (hidden ambiguity checks like FX direction/time basis),
  and must include clarification question plus fixed options.
- semantic_hints must be traceable to this case's function/question/context/article excerpt.
- Do not output generic options like "Option A/Option B".
- Keep semantic_hints concise and bounded: at most 6 hints per card.
- For each semantic_hint include:
  - id
  - ambiguity_type
  - trigger_signal
  - clarification_question
  - options  <- NOW STRUCTURED (see below)
  - i_level (hard|soft)
  - assumption_if_not_clarified
  - impact_scope (input_binding|output_interpretation|directionality)

Semantic hint options policy (IMPORTANT -- options are now objects, not strings):
Each option must be:
  { "label": "...", "value": "...", "is_default": true/false, "transform_spec": null | <spec> }

- Exactly one option must have is_default=true (the baseline the FIC code already implements).
- The DEFAULT option's transform_spec MUST be null. The FIC code IS the default -- no
  post-processing is needed for the baseline interpretation. Do not write a trivial
  identity transform like `result_expr: "result"` for the default.
- A non-default option MAY include a transform_spec ONLY IF the alternative interpretation
  has a pure-arithmetic relationship to the default output. If no such algebraic relationship
  exists (e.g. the alternative would require choosing a different input source, rescaling
  an INPUT before computation, or recomputing from scratch), the option MUST have
  transform_spec=null.
- result_expr is a PURE ARITHMETIC EXPRESSION. It MUST NOT contain function calls of any
  kind -- no compute(), no exp/log/sin/cos, no min/max, no abs. Only arithmetic operators
  (+, -, *, /, **) over `result` and input names. If the only way to express the
  alternative is via a function call, set transform_spec=null.
- Reference inputs by BARE NAME ONLY. Write `discount_rate`, not `inputs['discount_rate']`.
  Write `initial_investment`, not `inputs["initial_investment"]`. Subscript notation will
  be REJECTED by the verifier.
- ONLY use names that actually appear in this FIC's `inputs` list. If you need a quantity
  that is not in the inputs (e.g. `inflation_rate` for a real-rate transform on a FIC that
  only has nominal-rate inputs), set transform_spec=null -- do NOT fabricate input names.
- result_expr MUST reference `result`. It is a transformation OF the original output, not
  a recomputation of the formula. If the alternative requires recomputing from scratch
  (rather than adjusting `result`), set transform_spec=null.
- transform_spec schema:
  {
    "patch_type": "result_postprocess",
    "result_expr": "<Python expression using `result` and any input variable names>",
    "max_expr_nodes": <int, expected AST node count of result_expr>,
    "invariant": "<a Python `lhs == rhs` equality over result_old / result_new / input names>",
    "affected_inputs": ["<list of input names referenced in result_expr / invariant>"]
  }
- result_expr rules:
  * Use only `result` (the output of the original execution) and names that appear in the FIC inputs.
  * No function calls, no imports, no side effects. Pure arithmetic expression only.

WHEN TO USE transform_spec -- domain-driven guidance.
Pick the ambiguity that genuinely applies to THIS FIC's domain/inputs. Do NOT copy patterns
from unrelated domains. The hint must be traceable to the FIC's function/inputs.

| FIC domain / topic family                | Typical I-class ambiguity            | result_expr pattern                                  |
|------------------------------------------|--------------------------------------|------------------------------------------------------|
| time_value_money / annuity / NPV / PV    | end vs beginning of period (annuity) | `(result + initial_investment) * (1 + r) - initial_investment` |
| time_value_money / compounding           | discrete vs continuous compounding   | depends on formula; often `result * exp(r) / (1+r)`  |
| interest_rates / yield                   | nominal APR <-> effective APY          | `(1 + result/n) ** n - 1`                            |
| interest_rates / rate conversion         | monthly <-> annual rate                | `(1 + result) ** 12 - 1` or `(1 + result) ** (1/12) - 1` |
| any rate output                          | percent vs decimal scale             | `result / 100.0` or `result * 100.0`                 |
| fx / foreign_exchange                    | direct vs indirect quote             | `1.0 / result`                                       |
| cash_flow / valuation                    | inflow-positive vs outflow-positive  | `-result`                                            |
| returns / performance                    | pre-tax vs post-tax                  | `result * (1 - tax_rate)`                            |
| returns / performance                    | nominal vs real                      | `(1 + result) / (1 + inflation_rate) - 1`            |
| risk / volatility                        | period vol vs annualized             | `result * (252 ** 0.5)` or `result * (12 ** 0.5)`    |
| statistics / dispersion                  | sample vs population variance        | `result * (n - 1) / n`                               |
| bonds                                    | clean vs dirty price                 | `result + accrued_interest`                          |

If THIS FIC's domain is not in the table above, do NOT invent a forced transform_spec --
prefer transform_spec=null and let the system fall back to clarification-only.

ANTI-PATTERNS -- set transform_spec=null in these cases:
1. The ambiguity is about WHICH INPUT to use (e.g. "use WACC vs cost of equity" for
   discount_rate, "use S&P 500 vs Treasury" for benchmark). No post-processing can fix
   input choice -- the user must rerun with the right input.
2. The alternative interpretation requires re-running the formula with a structurally
   different input (e.g. "use only operating cash flows vs all cash flows").
3. The ambiguity is purely semantic categorization (e.g. "is this a hurdle rate or
   required return?") with no numerical adjustment.
4. The mapping would need to change loop logic or formula shape (not a closed-form
   post-process of `result`). Leave it null; future work will support code_patch.

WORKED EXAMPLE 1 -- NPV / annuity-due (time_value_money):
{
  "id": "payment_timing",
  "ambiguity_type": "time_basis",
  "trigger_signal": "payment period not explicitly stated as end-of-period",
  "clarification_question": "Are cash flows paid at END of each period (ordinary annuity) or BEGINNING (annuity-due)?",
  "i_level": "soft",
  "assumption_if_not_clarified": "End of period (ordinary annuity)",
  "impact_scope": "output_interpretation",
  "options": [
    { "label": "End of period -- DEFAULT", "value": "end_of_period", "is_default": true,  "transform_spec": null },
    {
      "label": "Beginning of period (annuity-due)", "value": "beginning_of_period", "is_default": false,
      "transform_spec": {
        "patch_type": "result_postprocess",
        "result_expr": "(result + initial_investment) * (1 + discount_rate) - initial_investment",
        "max_expr_nodes": 15,
        "invariant": "result_new + initial_investment == (result_old + initial_investment) * (1 + discount_rate)",
        "affected_inputs": ["discount_rate", "initial_investment"]
      }
    }
  ]
}

WORKED EXAMPLE 2 -- monthly <-> annual rate (interest_rates):
{
  "id": "rate_period_basis",
  "ambiguity_type": "rate_period",
  "trigger_signal": "rate may be quoted monthly or annually",
  "clarification_question": "Is the rate of return quoted on a monthly or annual basis?",
  "i_level": "soft",
  "assumption_if_not_clarified": "Annual rate",
  "impact_scope": "output_interpretation",
  "options": [
    { "label": "Annual -- DEFAULT", "value": "annual", "is_default": true, "transform_spec": null },
    {
      "label": "Monthly (annualize via compounding)", "value": "monthly", "is_default": false,
      "transform_spec": {
        "patch_type": "result_postprocess",
        "result_expr": "(1 + result) ** 12 - 1",
        "max_expr_nodes": 7,
        "invariant": "result_new == (1 + result_old) ** 12 - 1",
        "affected_inputs": []
      }
    }
  ]
}

WORKED EXAMPLE 3 -- percent vs decimal output (any rate-returning FIC):
{
  "id": "output_scale",
  "ambiguity_type": "unit_scale",
  "trigger_signal": "question asks for percentage, FIC outputs decimal (or vice versa)",
  "clarification_question": "Should the answer be expressed as a decimal (0.08) or percent (8.0)?",
  "i_level": "soft",
  "assumption_if_not_clarified": "Decimal",
  "impact_scope": "output_interpretation",
  "options": [
    { "label": "Decimal -- DEFAULT", "value": "decimal", "is_default": true, "transform_spec": null },
    {
      "label": "Percent", "value": "percent", "is_default": false,
      "transform_spec": {
        "patch_type": "result_postprocess",
        "result_expr": "result * 100.0",
        "max_expr_nodes": 3,
        "invariant": "result_new == result_old * 100.0",
        "affected_inputs": []
      }
    }
  ]
}

WORKED EXAMPLE 4 -- ANTI-PATTERN (benchmark choice, transform_spec MUST be null):
{
  "id": "benchmark_source",
  "ambiguity_type": "input_binding",
  "trigger_signal": "benchmark return source not explicitly named",
  "clarification_question": "Is the benchmark return based on the S&P 500, a sector ETF, or the risk-free rate?",
  "i_level": "soft",
  "assumption_if_not_clarified": "Market benchmark provided in inputs",
  "impact_scope": "input_binding",
  "options": [
    { "label": "Use given benchmark_return as-is -- DEFAULT", "value": "as_given", "is_default": true, "transform_spec": null },
    { "label": "Reconsider with alternative benchmark", "value": "alternative_benchmark", "is_default": false, "transform_spec": null }
  ]
}
(Why transform_spec is null here: choosing a different benchmark means re-running the formula
with different input values -- there is no closed-form expression mapping `result` from one
benchmark to another.)

Input case:
<DEFINITION_JSON>
<<SOURCE RECORD (JSON): source_meta, article_title, article_doc_id, article_content_excerpt, function, python_solution, context, question>>
</DEFINITION_JSON>

Taxonomy:
<TAXONOMY_JSON>
<<DOMAIN TAXONOMY (JSON)>>
</TAXONOMY_JSON>

Output:
Return JSON only.
\end{lstlisting}

\paragraph{Retrieval metadata (stage 2)}\leavevmode
\begin{lstlisting}
Role:
You are a retrieval engineer generating retrieval metadata for a financial formula card.

Task:
Given one `fic_core`, generate retrieval-focused fields only:
- title
- summary
- selection_hints
- applicable_when
- not_applicable_when
- scope_boundaries
- common_confusions
- keywords

Rules:
1) Optimize for card retrieval and disambiguation.
2) Keep statements concise and concrete.
3) Do not add execution rules or repair workflow details.
4) Keep topic family aligned with fic_core domain/topic.

Input:
<CORE_JSON>
<<CORE CARD FROM STAGE 1 (JSON)>>
</CORE_JSON>

Output:
Return JSON only.
\end{lstlisting}

\paragraph{Repair rules (stage 3)}\leavevmode
\begin{lstlisting}
Role:
You are a product workflow designer creating repair rules for diagnostic checks.

Task:
For each diagnostic check in `fic_core.diagnostic_checks`, generate a stable frontend repair rule.

Rules:
1) Each `rule_id` in diagnostic_checks must appear exactly once in output rules.
2) Keep actions deterministic and UI-friendly.
3) Focus on what to ask user and what next action to take.
4) Do not modify diagnostic_type/severity; they are inherited from core.
5) repair_action.type must be one of:
   request_missing_fields, confirm_assumption, confirm_unit_conversion,
   swap_suspected_fields, select_alternative_fic, rephrase_task_intent
6) allowed_next_steps must be from:
   rerun_same_fic, select_alternative_fic, ask_followup, stop_with_refusal

Input:
<CORE_JSON>
<<CORE CARD FROM STAGE 1 (JSON)>>
</CORE_JSON>

Output JSON shape:
{
  "rules": [
    {
      "rule_id": "...",
      "title": "...",
      "user_message": "...",
      "explanation": "...",
      "ask_user_for": [
        {"slot": "...", "label": "...", "type": "text|number|enum|boolean", "required": true, "options": [{"value":"...", "label":"..."}]}
      ],
      "repair_action": {"type": "...", "target": "..."},
      "allowed_next_steps": ["..."]
    }
  ]
}
\end{lstlisting}

\subsection{Authoring the human-labeled set}
GPT-5.2 wrote each query from a FinanceReasoning topic outside the 250, with one persona and one latent intent. The 18 queries of the missing-input stratum were later revised offline with a Codex agent, as described in Appendix~\ref{app:human}, and no model call was made for that revision.

\paragraph{Author prompt (GPT-5.2)}\leavevmode
\begin{lstlisting}
You are simulating a real person typing ONE financial question to an AI assistant.

Persona: <<PERSONA>>.

Topic to riff on (grounding only -- do NOT copy it or reuse its exact numbers):
  REFERENCE QUESTION: <<REFERENCE QUESTION>>
  REFERENCE CONTEXT: <<REFERENCE CONTEXT, FIRST 800 CHARACTERS>>

This time, the person's situation is: <<INTENT SITUATION>>
<<CONTEXT RULE>>

Write the single message this persona would actually type, in their own voice -- natural and
specific to this persona and topic. Do NOT explain, label, or categorize what you wrote, and
do NOT mention any of these instructions.

Return STRICT JSON only:
{"question": "<the user's message>", "context": "<numbers/background the user includes, or empty string>"}
\end{lstlisting}

\paragraph{Personas}\leavevmode
\begin{lstlisting}
hurried_analyst: a busy sell-side analyst firing off a quick question between meetings
retail_investor: a retail investor who is not sure of the exact financial terminology
ops_manager: a non-finance operations manager who just needs a number for a report
precise_quant: a precise quant analyst who states every input cleanly
headline_reader: someone who just read a news headline and wants to sanity-check one figure
\end{lstlisting}

\paragraph{Intents}\leavevmode
\begin{lstlisting}
clean
  situation: You have every figure you need and you ask clearly for one specific number.
  context rule: Put every number you use in `context`.

ambiguous
  situation: You give the numbers, but you are loose about which exact version you want: you say 'the return' / 'the growth' / 'the rate' / 'the margin' without pinning down annualized-vs-total, which basis (360 vs 365, gross vs net, per-share vs total), or which of several metrics (e.g. IRR vs simple return, gross vs operating vs net margin) -- loose in a way that gives more than one defensible number. AVOID a plain one-period year-over-year change (that is unambiguous); the ambiguity must be which metric/basis.
  context rule: Put the numbers in `context`, but keep the wording of WHAT you want genuinely under-specified so more than one answer is defensible.

judgment
  situation: You are NOT after a single computed number. You want an opinion or decision ('should I...', 'is it worth it', 'is this reasonable?'), a forecast ('will it...', 'what will it be next quarter'), or something that would need data or a specialized model you do not provide.
  context rule: Give whatever background you'd naturally mention in `context`; it need not be enough to compute anything.

missing (original; outputs revised, Appendix E)
  situation: You are looking at a report, table, or filing and you refer to it, but you do NOT include the actual figures in your message -- you only remember the topic, not the numbers, and you ask as if the assistant could see the document.
  context rule: Leave `context` EMPTY (or nearly so) -- you did not paste the numbers.
\end{lstlisting}
